\documentclass[11pt]{article}
\usepackage[T1]{fontenc}
\usepackage{textcomp}
\usepackage[margin=0.75in]{geometry}
\usepackage{amsmath,amssymb,amsthm}
\usepackage{array,booktabs}
\usepackage{hyperref}
\setlength{\parindent}{0pt}
\newtheorem{theorem}{Theorem}
\theoremstyle{definition}
\newtheorem{definition}{Definition}
\title{Flow Proof Artifact: Finset-inter-assoc}
\author{Generated by \texttt{flow doc proof}}
\date{Flow Proof Book}
\begin{document}
\maketitle
Finite set intersection associates.\\[0.5em]
\textbf{Source.} graham-knuth-patashnik — *Concrete Mathematics*\\[0.5em]
\bigskip
\noindent{\large \textbf{Derived fact 1.} \textit{(a ∩ b) ∩ c = a ∩ (b ∩ c)} \hfill Finset \cdot intersection \cdot parentheses do not matter}\par
\smallskip\noindent\textit{Coordinate: Finset \cdot intersection \cdot parentheses do not matter}\par
\smallskip\noindent\textit{Source: graham-knuth-patashnik}\par
\smallskip\noindent\textit{Built on: order does not matter, for intersection on Finset, empty is the left annihilator, for intersection on Finset}\par
\medskip
\noindent\textbf{Goal.} (a ∩ b) ∩ c = a ∩ (b ∩ c)\par
\noindent$\displaystyle \forall a \in Finset \forall b \in Finset \forall c \in Finset\quad (a ∩ b) ∩ c = a ∩ (b ∩ c)$\par
\medskip
\noindent
\begin{tabular}{@{} >{\raggedright\arraybackslash}p{0.06\textwidth} >{\raggedright\arraybackslash}p{0.47\textwidth} >{\raggedleft\arraybackslash}p{0.41\textwidth} @{}}
\textbf{\#} & \textbf{Proof} & \textbf{Mathematics} \\
\midrule
\textbf{1.} & We prove that parentheses do not matter for intersection on Finset. & \\[0.45em]
\textbf{2.} & We invoke the derived fact governing intersection on Finset: order does not matter, for intersection on Finset (instantiated for a, b). & \\[0.45em]
\textbf{3.} & We invoke the derived fact governing intersection on Finset: order does not matter, for intersection on Finset (instantiated for b, c). & \\[0.45em]
\textbf{4.} & We invoke the definitional clause governing intersection on Finset: empty is the left annihilator, for intersection on Finset (instantiated for c). & \\[0.45em]
\textbf{5.} & From step 2, step 3, and step 4, this implies (a ∩ b) ∩ c equals a ∩ (b ∩ c). Hence proven. & $\displaystyle (a ∩ b) ∩ c = a ∩ (b ∩ c)$ \\[0.45em]
\bottomrule
\end{tabular}
\label{thm:Finset-intersection-parentheses_do_not_matter}
\par\medskip\hrule
\end{document}
